% ============================================================
% CAPITOLO 14 - RIEPILOGO DELLA PARTE III
% ============================================================
\chapter{Riepilogo della Parte III}
\label{chap:riepilogo_iii}

\section{Sintesi degli usi del secondo rivolto}
\label{sec:sintesi_secondo_rivolto}

\begin{center}
\begin{tabularx}{\linewidth}{lLLLL}
\toprule
\textbf{Tipo} & \textbf{Posizione metrica} & \textbf{Movimento del basso} & \textbf{Ruolo} & \textbf{Dettaglio} \\
\midrule
Cadenziale & forte & fermo sulla dominante & prepara la dominante & Sez. \ref{sec:cadenziale} \\
Passaggio & debole & grado congiunto & collegamento melodico & Sez. \ref{sec:64_passaggio} \\
Pedale & indifferente & nota fissa & ornamento armonico & Sez. \ref{sec:64_pedale} \\
\bottomrule
\end{tabularx}
\end{center}

\section{Schema di funzioni, cadenze e frasi}
\label{sec:schema_iii}

\begin{center}
\small
\begin{tabular}{p{4.6cm}p{7.4cm}}
\toprule
\textbf{Argomento} & \textbf{Dove} \\
\midrule
Tonica, sottodominante, dominante & Sezione \ref{sec:tre_funzioni} \\
Sostituti delle funzioni principali & Sezioni \ref{sec:accordi_secondari} e \ref{subsec:accordi_funzioni} \\
Armonizzare la scala con le triadi & Sezione \ref{subsec:principio_regola_ottava} \\
Scelta dello stato dell'accordo & Regola \ref{reg:stato_basso} \\
Cadenze conclusive & Sezione \ref{sec:cadenze_conclusive} \\
Cadenze sospensive & Sezione \ref{sec:cadenze_sospensive} \\
Gerarchia delle cadenze & Regola \ref{reg:gerarchia_cadenze} \\
Modello T-S-D-T & Regola \ref{reg:tsdt} \\
Frase e periodo & Sezioni \ref{sec:frase_musicale} e \ref{sec:periodo} \\
Cadenze di metà e finale nel basso dato & Sezione \ref{sec:applicazione_frase} \\
\bottomrule
\end{tabular}
\end{center}
